% Part: incompleteness
% Chapter: introduction
% Section: overview

\documentclass[../../../include/open-logic-section]{subfiles}

\begin{document}

\olfileid{inc}{int}{ovr}

\olsection{Gegambaran Umum Asil Ora Jangkep}

Hilbert ngarepake yen matematika bisa diformalkan
ing teori kang !!{axiomatizable}, kang bisa
dibuktekake !!{complete} lan !!{decidable}.
Kajaba iku, dheweke arep mbuktekake konsistensi
teori iki nganggo cara kang ringkih banget,
yaiku cara ``finitèr'', kanggo mbela matematika
klasik saka tantangan intuisionisme. Teorema-teorema
G\"odel ngenani ora jangkep nuduhake yen
ancas-ancas iku ora bisa digayuh.

Teorema pisanan G\"odel ngenani ora jangkep
nuduhake yen salah sawijining versi
\emph{Principia Mathematica} anggitane Russell
lan Whitehead ora !!{complete}. Nanging buktine
sejatine umum banget lan bisa ditrapake marang
maneka warna teori. Iki ateges ora mung
\emph{Principia Mathematica} kang gagal nyakup
matematika kanthi pepak; \emph{ora ana} teori
kang konsisten, efektif diaksiomatisasi lan cukup
kuwat kanggo aritmetika kang bisa nindakake iku.
Perlu wektu kanggo nemtokake ciri-ciri teori
kang cukup supaya teorema-teorema ora jangkep
bisa ditrapake, lan kanggo nggeneralisasi bukti
G\"odel supaya mung gumantung marang ciri-ciri
iku. Saiki kita bisa nyatakake versi kang umum
banget saka teorema pisanan ngenani ora jangkep
kanggo teori ing basa aritmetika $\Lang L_A$.

\begin{thm}
Yen $\Gamma$ iku teori kang konsisten lan
!!{axiomatizable} ing~$\Lang L_A$ kang
!!{represents} kabeh fungsi komputabel lan
relasi !!{decidable}, mula $\Gamma$
ora !!{complete}.
\end{thm}

Ngandharake yen $\Gamma$ ora !!{complete} ateges
ana paling ora siji !!{sentence}~$!A$ saengga
$\Gamma \Proves/ !A$ lan $\Gamma \Proves/ \lnot
!A$. !!a{sentence} kaya mangkono diarani
\emph{mandiri} (saka~$\Gamma$).
Sejatine, ora angel banget mbuktekake yen ukara
mandiri mesthi ana. Nanging kakuwatan buktine
G\"odel yaiku menehi \emph{conto tartamtu}
saka !!{sentence} kang digawe kanthi cara efektif.
Konstruksi kang narik kawigaten iku ngasilake
!!a{sentence}~$!G_\Gamma$, kang diarani
\emph{ukara G\"odel} kanggo~$\Gamma$.
Ukara iki ora bisa dibuktekake amarga ing
$\Gamma$, $!G_\Gamma$ ekuivalen karo pratelan
yen $!G_\Gamma$ ora bisa dibuktekake ing~$\Gamma$. Cathetan koreksi sumber OLPL-775: konsistensi wae njamin ukara Gödel kang nyebut ora kabuktene dhewe ora bisa dibuktekake, nanging durung njamin negasine uga ora bisa dibuktekake. Kanggo kamandiran ukara wangun iki perlu syarat luwih kuwat, kayata konsistensi omega ing bukti asli; konstruksi Rosser bisa menehi ukara mandiri kanthi konsistensi wae. Prosa saiki ora maneh ngaku ukara Gödel biasa mesthi mandiri kanthi konsistensi wae.
Konstruksi iki ditindakake kanthi
\emph{konstruktif}: yen diwenehi aksiomatisasi
saka~$\Gamma$ lan katrangan sistem
!!{derivation}, bukti iku menehi cara nyata
kanggo nulis~$!G_\Gamma$.

Konstruksi ing buktine G\"odel mbutuhake cara
kanggo nyatakake ing $\Lang L_A$ sifat-sifat
lan operasi marang suku lan !!{formula}s saka
$\Lang L_A$ dhewe. Iki kalebu sifat kaya
``$!A$ iku !!a{sentence},'' ``$\delta$ iku
!!a{derivation} saka~$!A$,'' lan operasi kaya
$\Subst{!A}{t}{x}$. Cara iki kudu (a) nyatakake
sifat lan relasi kasebut liwat ``pangodean''
simbol lan urutane (yaiku wujud suku,
!!{formula}s, !!{derivation}s, lan liya-liyane)
minangka wilangan natural (kang bisa dirembug
dening $\Lang L_A$). Cara iku uga kudu (b)
nindakake pangodean kanthi cara saengga
$\Gamma$ bisa mbuktekake fakta kang cocog.
Mula kita kudu nuduhake yen sifat-sifat kasebut
dikodeake dening sifat wilangan natural kang
!!{decidable}, lan operasine cocog karo fungsi
komputabel ing wilangan natural. Iki diarani
``aritmetisasi sintaks''.

Nanging sadurunge nliti cara ngaritmetisasi
sintaks, kita bakal ngrembug syarat yen~$\Gamma$
``cukup kuwat'', yaiku !!{represents} kabeh
fungsi komputabel lan relasi !!{decidable}.
Iki mbutuhake definisi kang cetha tumrap
``komputabel''. Ana sawetara cara kanggo
netepake, umpamane liwat model mesin Turing
utawa minangka fungsi kang bisa diitung dening
program ing basa pamrograman umum. Nanging
amarga ancas kita yaiku supaya teori ing
basa~$\Lang L_A$ !!{represents}s fungsi
lan relasi kasebut, luwih becik milih
definisi komputabilitas kang prasaja kanggo
fungsi numeris wae. Iki gagasan
\emph{fungsi rekursif}. Mula kita bakal
ngrembug fungsi rekursif dhisik. Banjur kita
bakal nuduhake yen $\Th{Q}$ wis
!!{represents} kabeh fungsi lan relasi
rekursif. Iki ngidini kita ngetrapake
teorema ora jangkep marang teori tartamtu
kaya $\Th{Q}$ lan $\Th{PA}$, amarga kita
wis netepake yen teori-teori iki ``cukup kuwat''.

Asil pungkasan aritmetisasi sintaks yaiku
!!a{formula} $\OProv[\Gamma](x)$ kang,
liwat pangodean !!{formula}s minangka wilangan,
nyatakake kabuktenan saka aksioma~$\Gamma$.
Mligine, yen $!A$ dikodeake dening wilangan~$n$
lan $\Gamma \Proves !A$, mula
$\Gamma \Proves \Prov[\Gamma](\num{n})$.
``Predikat kabuktenan'' kanggo $\Gamma$ iki
uga ngidini kita nyatakake, ing teges tartamtu,
konsistensine $\Gamma$ minangka
!!a{sentence} ing~$\Lang L_A$:
tetepna ``ukara konsistensi'' kanggo~$\Gamma$
minangka !!{sentence}
$\lnot \Prov[\Gamma](\num{n})$, ing ngendi
$n$ dijupuk minangka kode sawijining
kontradiksi, umpamane~$\lfalse$. Teorema
kapindho ngenani ora jangkep nyatakake yen
teori-teori kang konsisten lan !!{axiomatizable}
kanthi syarat kabuktenan kang trep uga ora bisa
mbuktekake ukara konsistensine
dhewe. Syarat supaya teorema iki bisa
ditrapake rada luwih ketat tinimbang mung
teori kasebut nggambarake kabeh fungsi
komputabel lan relasi !!{decidable}, nanging
kita bakal nuduhake yen $\Th{PA}$ nyukupi.

\end{document}
